\paragraph{事件词的自由仿射单子与前缀仿射像的精确 Haar 体积}
在忠实的 Gödel--Tate 规范编码的前缀刚性与有限证书预算已经确定之后，事件词的连接本身还可在环境 Tate 模块上被直接重写为一个自由仿射子单子；与此同时，每个有限前缀所切出的环境仿射像也具有完全显式的 Haar 体积律。这样，符号连接、仿射作用与 \(2^L\)-进体积便在同一层面被统一为可直接计算的有限证书。

\begin{theorem}[事件词在 \texorpdfstring{$T_2(E_{\min})$}{T2(Emin)} 上生成自由仿射子单子]\label{thm:xi-time-part62cc-eventword-free-affine-submonoid}
沿用定理 \ref{thm:killo-godel-semidir-affine-2adic} 的记号。设 \(E\) 为有限事件集，
\[
\mathrm{code}:E\hookrightarrow \NN,
\qquad
M:=\max_{e\in E}\mathrm{code}(e),
\qquad
B:=2^L>M,
\]
并记
\[
T:=T_2(E_{\min})\cong \ZZ_2^2,
\qquad
v\in T\setminus 2T.
\]
对任意有限事件词
\[
a=a_1\cdots a_n\in E^\ast
\]
定义前缀向量与仿射映射
\[
x_a:=\sum_{t=1}^{n}\mathrm{code}(a_t)B^{t-1}v,
\qquad
F_a(x):=x_a+B^n x
\qquad(x\in T).
\]
并约定空词 \(\varnothing\) 满足 \(x_{\varnothing}=0\)、\(F_{\varnothing}=\mathrm{id}_T\)。再对无限事件流
\[
\mathbf e=(e_t)_{t\ge 1}\in E^{\NN}
\]
定义全息点
\[
X(\mathbf e):=\sum_{t\ge 1}\mathrm{code}(e_t)B^{t-1}v.
\]
则有：
\begin{enumerate}
  \item 对任意 \(a,b\in E^\ast\)，都有
  \[
  F_a\circ F_b=F_{ab},
  \]
  其中 \(ab\) 表示词连接；
  \item 对任意 \(a\in E^\ast\) 与 \(\mathbf e\in E^{\NN}\)，都有
  \[
  X(a\mathbf e)=F_a\bigl(X(\mathbf e)\bigr);
  \]
  \item 映射
  \[
  E^\ast\hookrightarrow \mathrm{Aff}(T),
  \qquad
  a\longmapsto F_a
  \]
  为单射。
\end{enumerate}
因此，\(\{F_a:a\in E^\ast\}\) 在 \(\mathrm{Aff}(T)\) 中构成一个与事件词自由单子 \(E^\ast\) 规范同构的自由仿射子单子。
\end{theorem}
\begin{proof}
设 \(a\in E^m\)、\(b\in E^n\)。由定义，
\[
x_{ab}
=
\sum_{t=1}^{m}\mathrm{code}(a_t)B^{t-1}v
+
\sum_{s=1}^{n}\mathrm{code}(b_s)B^{m+s-1}v
=
x_a+B^m x_b.
\]
故对任意 \(x\in T\) 都有
\[
F_a(F_b(x))
=
x_a+B^m(x_b+B^n x)
=
(x_a+B^m x_b)+B^{m+n}x
=
F_{ab}(x),
\]
从而得到第 \(1\) 条。

第 \(2\) 条同样直接计算即可：
\[
X(a\mathbf e)
=
\sum_{t=1}^{m}\mathrm{code}(a_t)B^{t-1}v
+
\sum_{s\ge 1}\mathrm{code}(e_s)B^{m+s-1}v
=
x_a+B^m X(\mathbf e)
=
F_a(X(\mathbf e)).
\]

最后证明单射。若 \(F_a=F_b\)，设 \(|a|=m\)、\(|b|=n\)。先取 \(x=0\)，得
\[
x_a=x_b.
\]
再对任意 \(x\in T\) 比较
\[
F_a(x)-F_a(0)=B^m x,
\qquad
F_b(x)-F_b(0)=B^n x,
\]
遂有
\[
B^m x=B^n x
\qquad(\forall x\in T).
\]
由于 \(T\cong \ZZ_2^2\) 无挠且 \(B>1\)，故必有 \(m=n\)。现设公共长度为 \(n\)。若 \(a\neq b\)，记 \(j\) 为首个满足 \(a_j\neq b_j\) 的位置，则
\[
x_a-x_b
=
B^{j-1}
\left(
\mathrm{code}(a_j)-\mathrm{code}(b_j)
+
\sum_{t>j}\bigl(\mathrm{code}(a_t)-\mathrm{code}(b_t)\bigr)B^{t-j}
\right)v.
\]
其中首项
\[
c_j:=\mathrm{code}(a_j)-\mathrm{code}(b_j)
\]
满足
\[
-(B-1)\le c_j\le B-1,
\qquad
c_j\neq 0,
\]
故 \(c_j\notin B\ZZ_2\)。而其余各项都落在 \(B\ZZ_2\) 中，因此括号内的 \(2\)-进系数不属于 \(B\ZZ_2\)。由 \(v\notin 2T\) 可知 \(v\notin BT\)，于是
\[
x_a-x_b\in B^{j-1}T\setminus B^jT,
\]
特别地 \(x_a-x_b\neq 0\)，与 \(x_a=x_b\) 矛盾。故只能有 \(a=b\)。单射性得证。证毕。
\end{proof}
\leanverified{paper\_xi\_time\_part62cc\_eventword\_free\_affine\_submonoid}

\begin{corollary}[前缀仿射像的精确 Haar 体积律]\label{cor:xi-time-part62cc-prefix-affine-haar-volume}
沿用定理 \ref{thm:xi-time-part62cc-eventword-free-affine-submonoid} 的记号。对任意长度 \(n\) 的事件词
\[
a\in E^n
\]
定义其前缀仿射像
\[
U_a:=F_a(T)=x_a+B^nT\subset T.
\]
若 \(m_T\) 为 \(T\cong \ZZ_2^2\) 上的归一化 Haar 概率测度，则：
\begin{enumerate}
  \item 每个 \(U_a\) 都是子群 \(B^nT\) 的一个陪集，因而亦即半径 \(B^{-n}\) 的闭球，并满足
  \[
  m_T(U_a)=B^{-2n};
  \]
  \item 对固定的 \(n\)，若 \(a\neq b\in E^n\)，则
  \[
  U_a\cap U_b=\varnothing;
  \]
  \item 因而长度 \(n\) 前缀的总体环境阴影满足
  \[
  m_T\!\left(\bigcup_{a\in E^n}U_a\right)=|E|^n B^{-2n}.
  \]
\end{enumerate}
\end{corollary}
\begin{proof}
由定义
\[
U_a=x_a+B^nT
\]
显然是 \(B^nT\) 的陪集。又 \(T\cong \ZZ_2^2\) 为秩 \(2\) 的自由 \(\ZZ_2\)-模，故
\[
T/B^nT\cong (\ZZ/2^{Ln}\ZZ)^2,
\]
从而
\[
|T:B^nT|=2^{2Ln}=B^{2n}.
\]
于是每个陪集的 Haar 质量都等于指数倒数 \(B^{-2n}\)，得到第 \(1\) 条。

再证第 \(2\) 条。若 \(U_a\cap U_b\neq\varnothing\)，则同为 \(B^nT\) 的陪集，必有
\[
U_a=U_b,
\qquad
x_a-x_b\in B^nT.
\]
若 \(a\neq b\)，记 \(j\le n\) 为其首个分歧位置。与上一定理证明中的同一分解完全相同，可得
\[
x_a-x_b\in B^{j-1}T\setminus B^jT.
\]
由于 \(j\le n\)，这尤其推出
\[
x_a-x_b\notin B^nT,
\]
与上式矛盾。故只能有 \(a=b\)，于是不同长度 \(n\) 词对应的 \(U_a\) 两两不交。

第 \(3\) 条于是由 \(|E|^n\) 个两两不交陪集的体积相加立即得到：
\[
m_T\!\left(\bigcup_{a\in E^n}U_a\right)
=
\sum_{a\in E^n}m_T(U_a)
=
|E|^n B^{-2n}.
\]
证毕。
\end{proof}

\noindent
因此，事件词的时间前缀在 Tate 模块中并非只是被动记录，而是以自由仿射单子的方式主动生成其环境几何；而每个长度 \(n\) 的前缀又恰好切出一个 Haar 质量为 \(B^{-2n}\) 的仿射球。于是，词连接的组合增长与 \(2^L\)-进环境体积之间便获得了完全显式的乘法对应。

\endinput
